\section{Reference Frame Rotation}
\label{SEC:Ref_Frame}

GFR supports a rotating reference frame in which all elements in the
domain move at a constant rotation rate.
%
This can be used to simulate flow within rotating channels or
turbomachinery components.
%
The default rotation rate is zero, which means the simulation is
computed in a fixed reference frame.

\begin{figure}[h]
 \begin{center}
 %\input{Figures/tikz_frames}
  \insertimage[width=0.75\textwidth]{tikz/tikz_frames}
  \caption{Reference Frames}
  \label{FIG:RefFrames}
 \end{center}
\end{figure}

%%%%%%%%%%%%%%%%%%%%%%%%%%%%%%%%%%%%%%%%%%%%%%%%%%%%%%%%%%%%%%%%%%%%%%%%%%%%%%%%

\Fref{FIG:RefFrames} illustrates the relationship between
the     inertial (fixed)    $(\fix{x},\fix{y},\fix{z})$ frame and
the non-inertial (rotating) $(\rot{x},\rot{y},\rot{z})$ frame,
whose origin is shifted by $\fvs{\FrameTransVecVar}$ and whose axes
may be rotated by the rotation rate vector $\fvs{\FrameSpinVecVar}$.
%
This spin vector $\fvs{\FrameSpinVecVar}$ is defined by
the rotation rate $\FrameRotRate$ of the reference frame
about a unit axis $\fvs{a}$.
\begin{align}
     \fvs{a}      =& \fvi{a}{x} \unitvec{\fix{x}}
                  +  \fvi{a}{y} \unitvec{\fix{y}}
                  +  \fvi{a}{z} \unitvec{\fix{z}}
   & \fvi{a}{}    =& \| \fvs{a} \| = \sqrt{\fix{a}_k \fix{a}_k} = 1
  \\
     \fvs{\FrameSpinVecVar} =& \fvi{\FrameSpinVecVar}{x} \unitvec{\fix{x}}
                            +  \fvi{\FrameSpinVecVar}{y} \unitvec{\fix{y}}
                            +  \fvi{\FrameSpinVecVar}{z} \unitvec{\fix{z}}
                            =  \fvi{\FrameSpinVecVar}{}  \fvs{a}
   & \fvi{\FrameSpinVecVar}{}
                     =&     \| \fvs{\FrameSpinVecVar}    \|
                     =  \sqrt{ \fvi{\FrameSpinVecVar}{k}
                               \fvi{\FrameSpinVecVar}{k} }
                     =  \FrameRotRate
\end{align}
%
A rotation rate tensor $\fts{\FrameSpinTenVar}$ can likewise be defined
such that
$\fts{\FrameSpinTenVar} \cdot \fvs{v}
   = \fvs{\FrameSpinVecVar} \crossprod \fvs{v}$
for any vector $\fvs{v}$.
%
For clarity, we shall write this using the following annotated tensor matrix
notation, where the associated basis vectors are indicated in blue to the left
and top of the rows and columns of each matrix.
%
\begin{align}
      \fti{\FrameSpinTenVar}{ij}
  =& - \perm{ijk} \fvi{\FrameSpinVecVar}{k}
  =  \begin{bNiceArray}{ccc}[first-row,first-col]
                          &  \unitvec{\fix{x}}
                          &  \unitvec{\fix{y}}
                          &  \unitvec{\fix{z}}   \\
       \unitvec{\fix{x}}  &  \mspace{3mu} 0
                          &           - \fvi{\FrameSpinVecVar}{\fix{z}}
                          &  \phantom{-}\fvi{\FrameSpinVecVar}{\fix{y}} \\
       \unitvec{\fix{y}}  &  \phantom{-}\fvi{\FrameSpinVecVar}{\fix{z}}
                          &  \mspace{3mu} 0
                          &           - \fvi{\FrameSpinVecVar}{\fix{x}} \\
       \unitvec{\fix{z}}  &           - \fvi{\FrameSpinVecVar}{\fix{y}}
                          &  \phantom{-}\fvi{\FrameSpinVecVar}{\fix{x}}
                          &  \mspace{3mu} 0
     \end{bNiceArray}
   & \fti{\FrameSpinTenVar}{}
                     =&     \| \fti{\FrameSpinTenVar}{ij}   \|
                     =  \sqrt{ \fti{\FrameSpinTenVar}{ij}
                               \fti{\FrameSpinTenVar}{ij}/2 }
                     =  \sqrt{ \fvi{\FrameSpinVecVar}{k}
                               \fvi{\FrameSpinVecVar}{k}    }
                     =  \FrameRotRate
\end{align}

%%%%%%%%%%%%%%%%%%%%%%%%%%%%%%%%%%%%%%%%%%%%%%%%%%%%%%%%%%%%%%%%%%%%%%%%%%%%%%%%

From this, a passive transformation $\fms{\FrameRotMatVar}$ can be constructed
for converting the vector components from the inertial to the
non-inertial unit vectors.
\begin{align}
     \label{EQN:QRotMatrixArbAxis}
     \fms{\FrameRotMatVar}
  =& \begin{bNiceArray}{ccc}[first-row,first-col]
                          &  \unitvec{\fix{x}}
                          &  \unitvec{\fix{y}}
                          &  \unitvec{\fix{z}} \\
       \unitvec{\rot{x}}  &  \fvia{x}^2        b +          c
                          &  \fvia{x} \fvia{y} b + \fvia{z} s
                          &  \fvia{x} \fvia{z} b - \fvia{y} s \\
       \unitvec{\rot{y}}  &  \fvia{x} \fvia{y} b - \fvia{z} s
                          &  \fvia{y}^2        b +          c
                          &  \fvia{y} \fvia{z} b + \fvia{x} s \\
       \unitvec{\rot{z}}  &  \fvia{x} \fvia{z} b + \fvia{y} s
                          &  \fvia{y} \fvia{z} b - \fvia{x} s
                          &  \fvia{z}^2        b +          c
     \end{bNiceArray}
  &
  &  \begin{aligned}
       \FrameRotAngle =& \int_{t=0}^{\, t} { \FrameRotRate dt } \\
       b =& 1 - \cos{ \FrameRotAngle } \\
       c =&     \cos{ \FrameRotAngle } \\
       s =&     \sin{ \FrameRotAngle }
     \end{aligned}
\end{align}
%
Note that $\fms{\FrameRotMatVar}(t)$ is a time-dependent proper orthogonal
rotational transformation tensor such that:
\begin{align}
    \label{EQN:QProperOrthoMatrix}
       {\fms{\FrameRotMatVar}}   {\fms{\FrameRotMatVar}}^T
 =& \; {\fms{\FrameRotMatVar}}^T {\fms{\FrameRotMatVar}} = {\bm I},
    \quad\quad
    \det {\fms{\FrameRotMatVar}} = 1
\end{align}
%
%% Next line is not that important.
%Because the determinant is one, it does not alter the magnitude of
%the vector on which it operates.

%%%%%%%%%%%%%%%%%%%%%%%%%%%%%%%%%%%%%%%%%%%%%%%%%%%%%%%%%%%%%%%%%%%%%%%%%%%%%%%%

The motion of this moving non-inertial reference frame $\rvs{x}$
can therefore be described in terms of the
arbitrary time-dependent translation and rotation of the
inertial reference frame $\fvs{x}$
by the following relations:
\begin{align}
    \rot{t}   =& \fix{t} + t_0 &
    \fix{t}   =& \rot{t} - t_0 \\
    \rvs{x}   =& {\fms{\FrameRotMatVar}}(t)
                 \left( \fvs{x} - \fvs{\FrameTransVecVar}(t) \right) &
    \fvs{x}   =& {\fms{\FrameRotMatVar}}^T
                        \rvs{x} + \fvs{\FrameTransVecVar}    \\
    \label{EQ:MovingCoords}
    \rot{x}_i =& \fmi{\FrameRotMatVar}{ij}
                 \left( \fvi{x}{j} - \fvi{\FrameTransVecVar}{j} \right) &
    \fix{x}_i =& \fmi{\FrameRotMatVar}{ij}^T
                        \rvi{x}{j} + \fvi{\FrameTransVecVar}{i}
\end{align}
%where $\rot{\phantom{.}}$ denotes quantites within the moving frame.

%%%%%%%%%%%%%%%%%%%%%%%%%%%%%%%%%%%%%%%%%%%%%%%%%%%%%%%%%%%%%%%%%%%%%%%%%%%%%%%%

It is important to remember that scalars, vectors, and tensors represent
physical quantities that are defined {\em with respect to a particular
reference frame} (\thatis as perceived by an observer in that frame).
%
For example, to an observer standing in a room with stagnant air, the air
is still. However, to another observer in the room riding a carousel, the
air would seem to approach at a speed equal and opposite his motion.
 
Once the physical orientation of a vector or tensor has been defined
(based on the chosen reference frame), it is still possible to represent
that quantity using a different set of unit vectors than the default set
for that frame.
%
This includes other unit vectors within the same frame
or even the unit vectors of another frame.
%
For example, using the transformation $\fms{\FrameRotMatVar}$,
it is possible to write any fixed frame vector or tensor
using the unit vectors of the moving frame.
\begin{align}
    \mvi{v}{i}  =& \fmi{\FrameRotMatVar}{im} \fvi{v}{m} \\
    \mti{T}{ij} =& \fmi{\FrameRotMatVar}{im} \fti{T}{mn}
                   \fmi{\FrameRotMatVar}{nj}^T
\end{align}
%
Thus, $\fvs{v}$ and $\mvs{v}$ represent the same physical vector
(magnitude and direction) but are written here with separate symbols
because their component values differ due to the different unit vectors.
%
An {\em objective} scalar, vector, or tensor is one that maintains the
same magnitude and orientation when observed from both the inertial and
non-inertial frames.  Such objective quantities transform according to:
\begin{align}
 && % empty column
    \label{EQN:ObjectiveScalar}
 &  \text{Scalar:} &
    \rot{s}     =& \fix{s}
 && % empty column
    \\
 && % empty column
    \label{EQN:ObjectiveVector}
 &  \text{Vector:} &
    \rvi{v}{i}  =& \mvi{v}{i}
 && % empty column
    \\ % = \fmi{\FrameRotMatVar}{im} \fvi{v}{m}  \\
 && % empty column
    \label{EQN:ObjectiveTensor}
 &  \text{Tensor:} &
    \rti{T}{ij} =& \mti{T}{ij}
%               =  \fmi{\FrameRotMatVar}{im} \fti{T}{mn}
%                  \fmi{\FrameRotMatVar}{nj}^T
 && % empty column
\end{align}
In other words, the vector/tensor defined with respect to the moving frame
perspective happens to match that observed from the fixed frame.
%
The same unit vectors must be used when comparing these vector/tensor component
values between frames, so the \emph{transformed} fixed frame components
must be used.

For these reasons, it is important to distinguish both the observer perspective
and the unit vectors used to express the component form of vectors and tensors.
%
To avoid confusion, we use the following convention:
quantities without an accent ($\fvs{v}$) are defined
with respect to the     inertial frame and that frame's unit vectors,
quantities denoted with ($\rvs{v}$) are defined
with respect to the non-inertial frame and that frame's unit vectors, and
quantities denoted with ($\mvs{v}$) are defined
with respect to the     inertial frame but written in terms of the
unit vectors in the non-inertial frame
through the application of the transformation tensor ${\fms{\FrameRotMatVar}}$.
%
Most of the scalars used herein are objective quantities,
and we shall not distinguish between such scalars in different frames.
%
However, because velocity is the most commonly encountered
non-objective vector quantity, this also means that
kinetic energy and any other quantity commputed from it
(\forexample total conditions for energy, enthalpy, pressure, or temperature)
are also non-objective scalar quantities.

%%%%%%%%%%%%%%%%%%%%%%%%%%%%%%%%%%%%%%%%%%%%%%%%%%%%%%%%%%%%%%%%%%%%%%%%%%%%%%%%

When applying the transformation to the rotation axis,
one finds (after substituting $\fms{\FrameRotMatVar}$ from
Eq.~(\ref{EQN:QRotMatrixArbAxis}) and doing some math):
\begin{align}
     \mvs{a}
  \equiv&
     \begin{bNiceArray}{c}[first-row,first-col]
                          &  1       \\
       \unitvec{\rot{x}}  &  \mvia{x}  \\
       \unitvec{\rot{y}}  &  \mvia{y}  \\
       \unitvec{\rot{z}}  &  \mvia{z}  \\
     \end{bNiceArray}
  =  {\fms{\FrameRotMatVar}} \fvs{a}
%%This is standard matrix notation.
% =  {\fms{\FrameRotMatVar}} \begin{bmatrix} a_x \\ a_y \\ a_z \end{bmatrix}
% =                          \begin{bmatrix} a_x \\ a_y \\ a_z \end{bmatrix}
%%This is annotated matrix notation.
  =  {\fms{\FrameRotMatVar}} \cdot
     \begin{bNiceArray}{c}[first-row,first-col]
                          &  1       \\
       \unitvec{\fix{x}}  &  \fvia{x}  \\
       \unitvec{\fix{y}}  &  \fvia{y}  \\
       \unitvec{\fix{z}}  &  \fvia{z}  \\
     \end{bNiceArray}
  =
     \begin{bNiceArray}{c}[first-row,first-col]
                          &  1       \\
       \unitvec{\rot{x}}  &  \fvia{x}  \\
       \unitvec{\rot{y}}  &  \fvia{y}  \\
       \unitvec{\rot{z}}  &  \fvia{z}  \\
     \end{bNiceArray}
     \label{EQN:RotationAxisRotation}
\end{align}
Thus, the values of the components of
$\mvs{a} = ( \mvia{x}, \mvia{y}, \mvia{z} )$
are identical to those of
$\fvs{a} = ( \fvia{x}, \fvia{y}, \fvia{z} )$
even though they are written using two different sets of unit vectors.
%
This occurs because as the frame rotates, the angle between the
rotation axis and the unit vectors remains the same.
%
Since $\fvs{\FrameSpinVecVar} = \FrameRotRate \fvs{a}$,
a similar conclusion holds for the vector components of
$\mvs{\FrameSpinVecVar}$ and $\fvs{\FrameSpinVecVar}$
as well as the tensor components of
$\mts{\FrameSpinTenVar}$ and $\fts{\FrameSpinTenVar}$.

\begin{align}
     \mvs{\FrameSpinVecVar}
  =& \fvi{\FrameSpinVecVar}{x} \unitvec{\rot{x}}
  +  \fvi{\FrameSpinVecVar}{y} \unitvec{\rot{y}}
  +  \fvi{\FrameSpinVecVar}{z} \unitvec{\rot{z}}
  =  \fms{\FrameRotMatVar}     \fvs{\FrameSpinVecVar}
  \\
      \mti{\FrameSpinTenVar}{ij}
  =& - \perm{ijk} \mvi{\FrameSpinVecVar}{k}
  =  \begin{bNiceArray}{ccc}[first-row,first-col]
                          &  \unitvec{\rot{x}}
                          &  \unitvec{\rot{y}}
                          &  \unitvec{\rot{z}}   \\
       \unitvec{\rot{x}}  &  \mspace{3mu} 0
                          &           - \fvi{\FrameSpinVecVar}{\fix{z}}
                          &  \phantom{-}\fvi{\FrameSpinVecVar}{\fix{y}} \\
       \unitvec{\rot{x}}  &  \phantom{-}\fvi{\FrameSpinVecVar}{\fix{z}}
                          &  \mspace{3mu} 0
                          &           - \fvi{\FrameSpinVecVar}{\fix{x}} \\
       \unitvec{\rot{x}}  &           - \fvi{\FrameSpinVecVar}{\fix{y}}
                          &  \phantom{-}\fvi{\FrameSpinVecVar}{\fix{x}}
                          &  \mspace{3mu} 0
     \end{bNiceArray}
  =  ( \fms{\FrameRotMatVar} \fts{\FrameSpinTenVar} {\fms{\FrameRotMatVar}}^T )_{ij}
\end{align}

%%%%%%%%%%%%%%%%%%%%%%%%%%%%%%%%%%%%%%%%%%%%%%%%%%%%%%%%%%%%%%%%%%%%%%%%%%%%%%%%

For additional information about the derivation of the equations of
motion in a rotating reference frame, see \linecite{Yoder_2026}.

%%%%%%%%%%%%%%%%%%%%%%%%%%%%%%%%%%%%%%%%%%%%%%%%%%%%%%%%%%%%%%%%%%%%%%%%%%%%%%%%

\SeeAlso{
\inp[nmlPhysics]{nml\_physics},
\inp[nmlPhysics::MovingRefFrame/Rotation]%
    {Moving\_Ref\_Frame\%Rotation},
\inp[nmlPhysics::MovingRefFrame/Translation]%
    {Moving\_Ref\_Frame\%Translation}
}

